% =====================================================================
% Interpretability as a Science (NeurIPS 2026, Sydney) -- SHORT PAPER
% Limit: 5 pages excluding references and appendices. Non-archival.
% Submission: OpenReview, NeurIPS.cc/2026/Workshop/InterpScience
% Deadline: 2026-08-28 AoE
%
% TO SWITCH TO THE OFFICIAL TEMPLATE: download neurips_2026.sty, then
% replace the PREAMBLE block below with:
%     \documentclass{article}
%     \usepackage{neurips_2026}          % anonymous by default
% The body below needs no other changes.
% =====================================================================

\documentclass[10pt]{article}
\usepackage[margin=1in]{geometry}
\usepackage{times}
\usepackage{booktabs}
\usepackage{amsmath,amssymb}
\usepackage{graphicx}
\usepackage{xcolor}
\usepackage[colorlinks=true,allcolors=blue]{hyperref}
\usepackage{natbib}
\usepackage{caption}
\captionsetup{font=small,labelfont=bf}
\setlength{\parskip}{2pt}

\title{Semantically Valid, Behaviorally Predictive, Causally Inert:\\
A Validity Ladder for Emotion Probes in Llama 3.1 70B}

\author{Anonymous Author(s)\\Affiliation withheld for double-blind review}
\date{}

\begin{document}
\maketitle

\begin{abstract}
Interpretability results are routinely reported as though semantic
interpretability, behavioral prediction, and causal control were a single
property of a learned direction. We present a case study in which they come
apart. Extracting linear emotion representations from the residual stream of
Llama~3.1~70B and studying reward hacking on a code-generation task, we
evaluate the same set of probes against four progressively stricter validity
criteria. The probes recover human valence judgments on an external corpus
(cross-validated $R^2 = 0.38$, $n = 10{,}062$). Their natural, unsteered
activation predicts subsequent reward hacking with AUC $0.832$
$[0.742, 0.909]$ (permutation $p < 0.001$), a result that reproduces exactly
across two independent judge models from different families ($80/80$ trial-level
agreement). Yet steering along the same directions does not change behavior more
than steering along a random direction of equal norm ($7.5\%$ vs.\ $6.5\%$ hack
rate; Fisher $p = 0.835$), and the fine-grained dose--response is flat
($z = -1.42$, $p = 0.157$). The incremental advantage over a chain-of-thought
text baseline, meanwhile, is strongly layer-dependent and is not significant at
our pre-specified layer. We report the trajectory by which our own initial
claims degraded under each added control, and argue that semantic, predictive,
and causal validity should be reported as separate, separately-evidenced tiers
rather than inferred from one another.
\end{abstract}

\section{Introduction}

A common inferential shortcut in interpretability runs: we found a direction that
is semantically interpretable; it predicts a behavior we care about; therefore it
is a mechanism, and steering it is a lever. Each link in that chain is an
empirical claim, and each can fail independently.

\citet{sofroniew2026} report that Claude Sonnet 4.5 forms linear representations
of emotion concepts that causally drive alignment-relevant behavior, with
steering along a ``desperate'' direction increasing reward hacking, and note that
such behavioral shifts can occur without visible trace in the model's output
text. This is an unusually clean setting in which to ask the follow-up question:
\emph{if we replicate the representation, do we also replicate the causal
handle?}

We extract emotion vectors from the open-weight Llama~3.1~70B and evaluate them
against four tiers of validity (Table~\ref{tab:ladder}). Our contribution is not
a new method. It is a documented dissociation: on the same probes, the same
trials, and the same behavioral labels, tiers 1 and 2 hold, tier 3 holds only
under a post-hoc layer choice, and tier 4 fails outright. We also report the
sequence in which our own earlier, stronger conclusions collapsed --- which
control killed which claim --- because that sequence is more diagnostic than the
endpoint alone.

\begin{table}[t]
\centering
\small
\begin{tabular}{@{}clll@{}}
\toprule
& \textbf{Tier} & \textbf{Question} & \textbf{Verdict} \\
\midrule
1 & Semantic     & Do probes track what humans call emotion? & Holds \\
2 & Predictive   & Does natural probe state predict the behavior? & Holds \\
3 & Incremental  & Does it beat a text-based baseline? & Layer-dependent \\
4 & Causal       & Does steering it beat a random direction? & \textbf{Fails} \\
\bottomrule
\end{tabular}
\caption{The validity ladder. Each tier is evidenced separately in \S\ref{sec:results}.}
\label{tab:ladder}
\end{table}

\section{Setup}

\paragraph{Probes.} We extract 50 emotion vectors from the Llama~3.1~70B residual
stream by contrasting emotion-eliciting narrative contexts, with PCA denoising.
Our pre-specified probe site is layer 53 of 80. $V_{\text{internal}}$ for a trial
is the cosine similarity between the trial's residual-stream activation and a
frozen Phase-1 emotion vector. \textbf{Probes are never refit on behavioral
data}, so every predictive result below is a transfer claim rather than a
within-dataset classification score.

\paragraph{Behavior.} Reward hacking is elicited with impossible-by-construction
coding tasks, where an honest attempt must fail and a shortcut (hardcoding
expected outputs, faking a verifier, silently dropping a requirement) passes.
A parallel sycophancy task was collected but is excluded throughout: it yielded
2 sycophantic responses in 650 trials and has no usable variance.

\paragraph{Labels.} Responses are classified by an LLM judge with three
independent passes. Because judge choice turned out to be consequential, we run
\emph{two} judges from different model families --- Qwen~2.5~72B and Claude ---
and report their agreement wherever labels carry a conclusion.

\paragraph{Text baseline.} $V_{\text{text}}$ is an 8-dimension emotional tone
rating (1--7) of the model's chain-of-thought, produced by the same judge and
blind to steering condition. It stands in for what a CoT monitor could observe.

\paragraph{Preregistration.} The tier-2 hypothesis was preregistered with a
binding decision rule before the confirmatory data were collected. On the
original dataset the rule returned \textsc{insufficient-data} (only 2 of 4 task
variants produced outcome variance); we honored that verdict and collected 80
additional unsteered trials rather than reporting the exploratory value.

\section{Results}
\label{sec:results}

\subsection{Tier 1: the probes are semantically valid}

We project the frozen probes onto EmoBank, a corpus of 10,062 sentences with
human valence--arousal--dominance ratings, with no refitting.

\begin{center}\small
\begin{tabular}{@{}lccc@{}}
\toprule
& Valence & Arousal & Dominance \\
\midrule
Cross-validated $R^2$ & \textbf{0.377} & 0.180 & 0.154 \\
\bottomrule
\end{tabular}
\end{center}

Valence is recovered substantially; arousal and dominance weakly. The sign
structure is coherent (\texttt{ecstatic} $r = +0.38$, \texttt{excited} $+0.36$,
\texttt{furious} $-0.34$, \texttt{terrified} $-0.34$). The directions are
measuring something a human would recognize as affect, not an arbitrary
task-correlated axis --- at least along valence.

\subsection{Tier 2: natural probe state predicts reward hacking}

On unsteered trials only ($n = 120$, 14 hacking events), with no intervention of
any kind:

\begin{center}\small
\begin{tabular}{@{}lcccc@{}}
\toprule
Predictor & AUC & Bootstrap 95\% CI & Permutation $p$ & Mann--Whitney $p$ \\
\midrule
$V_{\text{internal}}$[desperate] & \textbf{0.832} & $[0.742,\ 0.909]$ & $< 0.001$ & $2.8\times10^{-5}$ \\
$V_{\text{internal}}$[$-$calm]   & 0.865 & --- & --- & --- \\
\bottomrule
\end{tabular}
\end{center}

This is the paper's one confirmatory positive. It is robust to judge identity:
Qwen and Claude agree on \textbf{80/80 trials (100\%)} on this set, flagging the
same shortcuts. A representation extracted with no access to the behavioral task
predicts, from the residual stream, whether the model is about to cheat.

\subsection{Tier 3: incremental validity over CoT text is layer-dependent}

Pooled across steered and unsteered trials ($n = 992$, 87 events),
$V_{\text{internal}}$ substantially outperforms the text baseline: AUC $0.836$
$[0.797, 0.872]$ vs.\ $0.647$ $[0.603, 0.690]$, paired
$\Delta = 0.189$ $[0.141, 0.239]$. The safety-relevant partition is stark: among
453 trials where the probe reads high but the text reads calm, the hack rate is
$14.6\%$, against $0.8\%$ where both read low.

This pooled comparison, however, inherits a task-identity confound (probes
predict which task is being solved at $95.4\%$ accuracy; residualizing on task id
reduces AUC from $0.799$ to $0.674$). Restricting to the clean unsteered set and
sweeping the probe layer (Table~\ref{tab:layers}) gives a less comfortable
picture.

\begin{table}[t]
\centering\small
\begin{tabular}{@{}lcccc@{}}
\toprule
Layer & $V_{\text{int}}$[desp] AUC & $\Delta$ vs.\ $V_{\text{text}}$ & 95\% CI & $p$ (1-sided) \\
\midrule
13 & 0.691 & $+0.053$ & $[-0.173,\ 0.299]$ & 0.339 \\
26 & 0.450 & $-0.187$ & $[-0.493,\ 0.125]$ & 0.890 \\
\textbf{39} & \textbf{0.918} & $\mathbf{+0.280}$ & $\mathbf{[0.056,\ 0.479]}$ & \textbf{0.008} \\
52 & 0.708 & $+0.069$ & $[-0.247,\ 0.343]$ & 0.306 \\
53 (pre-specified) & 0.677 & $+0.037$ & $[-0.282,\ 0.329]$ & 0.391 \\
65 & 0.472 & $-0.165$ & $[-0.539,\ 0.202]$ & 0.809 \\
\bottomrule
\end{tabular}
\caption{Incremental AUC over the text baseline on unsteered trials
($n = 80$, 7 events). \textbf{At our pre-specified layer 53 the gap is not
significant.} Layer 39 clears $p < 0.01$, but with 6 layers tested the Bonferroni
threshold is $\alpha = 0.0083$: it passes under Qwen labels ($p = 0.0080$) and
fails under Claude labels ($p = 0.0097$).}
\label{tab:layers}
\end{table}

We therefore decline to claim a confirmed faithfulness gap. The honest statement
is: \emph{the gap is large when trials are pooled, absent at the pre-specified
layer on clean trials, and present at one exploratory layer at a significance
level that does not survive correction robustly.} Layer 39 is a hypothesis for
replication, not a result.

\subsection{Tier 4: causal specificity fails}

The decisive control. We steer along five random directions orthogonal to the
emotion subspace, at the same norm as the emotion steering.

\begin{center}\small
\begin{tabular}{@{}lcc@{}}
\toprule
Condition & hacks / $n$ & rate \\
\midrule
Unsteered baseline & 14 / 120 & 11.7\% \\
Emotion vector @ $\pm0.3$ & 12 / 160 & 7.5\% \\
Random directions @ $\pm0.3$ & 13 / 200 & 6.5\% \\
\midrule
\multicolumn{3}{@{}l@{}}{\emph{Emotion vs.\ random: Fisher exact $p = 0.835$}} \\
\bottomrule
\end{tabular}
\end{center}

Steering the emotion direction does nothing that steering an arbitrary direction
of the same magnitude does not also do. The fine-grained dose--response over
$\pm0.05$ to $\pm0.5$ is flat for the desperate direction ($z = -1.42$,
$p = 0.157$); the calm direction reaches $p = 0.041$ uncorrected with the wrong
sign and does not survive correction. Injecting ``feel desperate'' into the
system prompt likewise fails to separate from baseline ($17.5\%$ vs.\ $11.7\%$,
Fisher $p = 0.417$).

We emphasize what this does \emph{not} say. It does not show that emotion
representations are epiphenomenal in general, and it does not contradict
\citet{sofroniew2026} on a different model. It shows that in this model, at this
site, with this norm, \textbf{a direction that is semantically valid and
behaviorally predictive supplies no causal handle} --- and that without a
random-direction control we would have reported the opposite.

\subsection{What died, and to what}

Our own earlier analysis of this dataset reported a cross-validated AUC of
$0.997$, a leave-one-task-out AUC of $0.992$, and a causal steering effect at
$p = 0.005$. Table~\ref{tab:autopsy} records which control removed which claim.

\begin{table}[t]
\centering\small
\begin{tabular}{@{}p{0.35\textwidth}p{0.28\textwidth}p{0.28\textwidth}@{}}
\toprule
\textbf{Original claim} & \textbf{Control applied} & \textbf{Outcome} \\
\midrule
CV AUC $0.997$, LOTO $0.992$ & Cross-mechanism transfer & AUC $0.303$, below chance \\
Steering effect $p = 0.005$ & BH-FDR over $40+$ tests & $p = 0.11$, n.s. \\
Steering direction reversal is a novel model-specific finding & Random-direction null & Indistinguishable, $p = 0.835$ \\
Within-task AUC $> 0.98$ proves the signal & Task-id residualization & $0.799 \rightarrow 0.674$ \\
H5 confirmed at AUC $0.955$ & Preregistered decision rule & \textsc{insufficient-data}; re-run after new collection \\
\bottomrule
\end{tabular}
\caption{Claim autopsy. No claim survived every control it was subjected to in
its original form; one (tier 2) survived in weakened form after additional data
collection.}
\label{tab:autopsy}
\end{table}

\section{Discussion}

\paragraph{Prediction is not control, and interpretability is not either.}
The tier-2 result is genuinely useful for monitoring: a probe that flags
impending reward hacking at AUC $0.83$ with no intervention is a plausible
detector. The tier-4 failure means the same probe is not a plausible
\emph{intervention} target. These are different products and should be evidenced
differently. We suspect the conflation is common because the same figure ---
a direction with a name --- serves both narratives.

\paragraph{Random-direction controls should be default, not optional.}
Our causal claim was significant, replicated a published effect qualitatively,
and was wrong. The only thing that revealed it was steering along directions with
no semantic content at matched norm. Any steering result without such a control
is compatible with ``large perturbations perturb.''

\paragraph{Baselines need measured resolution.}
Part of our tier-3 gap is an artifact of the baseline: our $V_{\text{text}}$
distribution has 946 trials in the low third, \emph{zero} in the middle third,
and 46 in the high third, across only 14 unique rounded values. A comparison
against a near-degenerate baseline overstates the gap. We report the instrument's
resolution alongside the comparison; we think this should be routine.

\paragraph{Judge agreement is a load-bearing statistic.}
On the unsteered set, two judges from different families agreed on $80/80$
trials, which is why we trust tier 2. On a second dataset of 650 trials designed
to test cross-mechanism generalization, the same two judges agreed on
$165/650$ ($25.4\%$): Qwen labeled 2 responses as shortcuts and 524 as unclear,
while Claude labeled 194 as shortcuts. We consequently withhold \emph{all}
generalization conclusions from that dataset, including the negative one, and
recommend reporting cross-family judge agreement per-dataset rather than per-paper.

\section{Limitations}

Single model, single behavior, single probe-extraction method; 14 events at
$n = 120$ for the headline result. The tier-4 null is a null: it bounds the
effect at this site and norm but cannot exclude effects at other layers, larger
norms, or with steering applied at other token positions. Layer 39 is a post-hoc
selection from a 6-layer sweep and is reported as such. The cross-mechanism
dataset is unusable pending judge reconciliation, so we cannot say whether the
tier-2 result generalizes across shortcut types --- in our view the single most
important open question, and one we deliberately do not answer here rather than
answer from unreliable labels.

\section{Conclusion}

We report a clean dissociation between semantic, predictive, and causal validity
for the same set of interpretability probes, together with the sequence of
controls that produced it. We suggest that these tiers be reported separately and
evidenced separately, that random-direction nulls accompany steering claims, that
baseline instruments have their resolution stated, and that cross-family judge
agreement be reported per dataset. Our strongest positive result --- residual-stream
state predicting reward hacking before it occurs --- survives all of this; our
most exciting one did not.

\paragraph{Reproducibility.} All analysis scripts, judged trial records,
preregistration, and a claim-to-file provenance table are released with the paper.

\bibliographystyle{plainnat}
\begin{thebibliography}{1}
\bibitem[Sofroniew et~al.(2026)]{sofroniew2026}
N.~Sofroniew, I.~Kauvar, et~al.
\newblock Emotion concepts and their function in a large language model.
\newblock \emph{Transformer Circuits Thread}, 2026.
\end{thebibliography}

\end{document}
